\chapter{Intrathecal Baclofen Pump Therapy}

\section*{Overview}
Intrathecal baclofen therapy delivers the muscle relaxant baclofen directly into cerebrospinal fluid through an implanted pump and catheter.

\section*{Why This Test or Procedure Is Ordered}
It is used for severe spasticity that is inadequately controlled by oral medicines or causes functional disability, painful spasms, or difficult caregiving.

\section*{How This Test or Procedure Is Performed}
After a screening dose or trial, a programmable pump is implanted under the abdominal skin and connected to a catheter entering the spinal canal. The dose is adjusted electronically and the pump is refilled periodically.

\section*{Risks and Follow-up}
Infection, catheter malfunction, overdose, withdrawal, respiratory depression, and need for revision are possible. Abrupt interruption can cause life-threatening withdrawal, so refills and emergency planning are essential.

\section*{References}
\begin{enumerate}
  \item MedlinePlus. Intrathecal Baclofen. U.S. National Library of Medicine; 2025.
  \item Current Medical Diagnosis and Treatment. McGraw-Hill; 2025.
\end{enumerate}
\clearpage
